\documentclass[11pt]{article}

\usepackage[margin=1in]{geometry}
\usepackage[T1]{fontenc}
\usepackage[utf8]{inputenc}
\usepackage{lmodern}
\usepackage{microtype}
\usepackage{amsmath, amssymb, amsthm}
\usepackage{booktabs}
\usepackage{array}
\usepackage{graphicx}
\usepackage{xcolor}
\usepackage{hyperref}
\usepackage{enumitem}
\usepackage{caption}
\usepackage{subcaption}
\usepackage{longtable}
\usepackage{float}
\usepackage{adjustbox}

\hypersetup{
    colorlinks=true,
    linkcolor=blue!50!black,
    citecolor=blue!50!black,
    urlcolor=blue!50!black
}

\setlist[itemize]{leftmargin=1.5em}
\setlist[enumerate]{leftmargin=1.5em}

\newcommand{\hphase}{h_1(T)}
\newcommand{\modsector}{s_{6,M}}
\newcommand{\zprimary}{2.575}
\newcommand{\zsensitivity}{2.326}
\newcommand{\Arc}{\operatorname{Arc80}}
\newcommand{\R}{\mathcal{R}}
\newcommand{\Ssector}{S_{\mathrm{sector}}}
\newcommand{\Sint}{S_{\mathrm{int}}}
\newcommand{\Smain}{S_{\mathrm{main}}}
\newcommand{\principaldoi}{10.5281/zenodo.19827222}
\newcommand{\paperdoi}{10.5281/zenodo.20017507}
\newcommand{\terminalsha}{c80b841b1ebd0a602b45db9fa0544499fc6781858b70bbb1ec0940aeffe0fa76}

\newcommand{\optionalfigure}[3]{%
\begin{figure}[H]
    \centering
    \IfFileExists{#1}{%
        \includegraphics[width=0.98\textwidth]{#1}%
    }{%
        \fbox{\parbox[c][2.25in][c]{0.88\textwidth}{\centering
            Figure placeholder: \texttt{#1}}}%
    }
    \caption{#2}
    \label{#3}
\end{figure}
}

\title{A Search-Adjusted Audit of Row-Sector Modulus Signals in Terminal Zeta-Zero Spacing Residuals}
\author{Jeffery Huckstead\\Independent Researcher}
\date{Replacement preprint draft}

\begin{document}

\maketitle

\begin{abstract}
We re-examine a previously reported row-sector structure in terminal adjacent-zero spacing residuals of the Riemann zeta function. The earlier preprint identified an apparent base-6 row-sector rebound near a \(\sim 10^3\)-row seam and reported modulus \(1009\) as the leading sector-amplitude ruler, with modulus \(990\) co-significant under local permutation guardrails. Subsequent validation showed that the principal companion paper's apparent \(h_1\) phase-locking structure is better explained as a projection of terminal altitude localization. This motivated the corresponding question for the row-sector claim: does the reported modulus-\(1009\) structure survive a look-elsewhere null that searches the same modulus family?

We answer this question negatively. A dense validation audit scans all moduli \(M=800,\ldots,1200\) and computes the sector-amplitude statistic \(S_{\mathrm{sector}}(M)\). Each permutation null shuffles tail labels and is allowed to scan the same 401 moduli, retaining its maximum sector amplitude. Under this search-adjusted null, neither modulus \(1009\) nor the best real modulus survives in any tested seam-scale configuration. In four primary fitted-ratio windows \(W\in\{900,1000,1100,1300\}\), modulus \(1009\) survives zero times, the best real modulus survives zero times, and modulus \(1009\) is never ranked in the top ten.

The replacement conclusion is therefore methodological rather than positive: the earlier row-sector signal is not a search-hardened discovery. It is consistent with a look-elsewhere artifact arising from sparse tail-sign populations searched over many moduli. Together with the revised companion result, this paper supports a unified interpretation: the surviving empirical object is terminal altitude architecture in highly compressed adjacent-spacing packets, while both continuous \(h_1\) phase structure and discrete modulus-\(1009\) sector structure were projection/search artifacts generated by exploratory coordinates.
\end{abstract}

\noindent\textbf{Keywords:} Riemann zeta zeros; zero spacing; row-sector residuals; modulus search; look-elsewhere effect; post-selection inference; terminal band; empirical audit.

%--------------------------------------------------------------------
\section{Introduction}
%--------------------------------------------------------------------

The present paper replaces an earlier positive interpretation of row-sector rebound dynamics in terminal zeta-zero spacing residuals. The original version reported a discrete row-index sector structure, amplitude-dominant under modulus \(1009\), co-significant under modulus \(990\), and empirically orthogonal to a previously reported \(h_1\) packet organization. That interpretation is no longer maintained.

The change follows a blind-review objection concerning post-selection inference. The central issue is not whether a chosen modulus can look structured after being selected, but whether the structure survives a null model that is allowed to search the same modulus family. A local permutation null conditional on a selected modulus does not address the look-elsewhere effect created by scanning many possible moduli, window sizes, metrics, and sectorizations.

This paper therefore asks a narrower and harder question:
\begin{quote}
Does the apparent modulus-\(1009\) row-sector signal survive a dense search-adjusted null over \(M=800,\ldots,1200\)?
\end{quote}

The answer is no. Once the null is allowed to scan the same modulus range and keep its own best result, the real modulus-\(1009\) signal is fully absorbed by the search-adjusted null envelope.

\subsection{Relation to the revised companion paper}

The companion paper originally framed rare extreme compressions as \(h_1\)-phase-locked. Subsequent validation audits rejected that interpretation. The compressed packet population itself survived, but the \(h_1\) corridor was explained as a projection of terminal altitude localization. The revised companion claim is therefore not ``\(h_1\) phase-locking,'' but rather:
\begin{quote}
highly compressed adjacent-spacing packets form a finite-horizon terminal altitude architecture.
\end{quote}

This paper performs the corresponding cleanup for the row-sector continuation. Its conclusion is stronger as a negative result than as the earlier positive claim: the row-sector modulus signal does not survive search-adjusted modulus controls and should not be presented as a discovered sector engine.

\subsection{Scope}

All statements are empirical and finite-horizon. No asymptotic theorem is claimed. The analysis concerns the loaded terminal LMFDB-derived board spanning approximately \(28.499547\mathrm{B}\) to \(30.610043\mathrm{B}\). The filename suffix \texttt{31p5B} used in computational artifacts denotes a downloader ceiling, not a source-data endpoint.

%--------------------------------------------------------------------
\section{Data and terminal board}
%--------------------------------------------------------------------

\subsection{Source board}

All computations use the seam-stitched terminal checkpoint parquet,
\begin{center}
\texttt{df\_master\_terminal\_checkpoint\_v16c\_8000\_31p5B.parquet},
\end{center}
containing \(935{,}111\) terminal rows with altitude range approximately
\[
28.499547\mathrm{B}\le T\le 30.610043\mathrm{B}.
\]
The seam-stitching preserves adjacent-gap continuity across source shard boundaries. The suffix \texttt{31p5B} is a retrieval ceiling only. A closure audit recorded SHA-256 checksum
\[
\texttt{\terminalsha}.
\]

\subsection{Residual coordinates}

The row-window residual pipeline uses spacing-ratio coordinates inherited from the companion audit chain:
\begin{align}
    \mathrm{Gap\_to\_Pure\_Ratio}   &= \frac{\Delta_k}{s_{\mathrm{pure}}(T_k)},\\
    \mathrm{Gap\_to\_Fitted\_Ratio} &= \frac{\Delta_k}{\widehat{s}(T_k)}.
\end{align}
The primary dense-modulus validation uses \(\mathrm{Gap\_to\_Fitted\_Ratio}\), which was the principal fitted-spine metric in the validation chain.

For a window of size \(W\), row-window means are formed, a slow altitude trend is removed by a fitted log-altitude baseline, and residuals are robustly standardized. The tail population consists of rows satisfying
\[
|z|\ge \zprimary.
\]

%--------------------------------------------------------------------
\section{Original row-sector instrument}
%--------------------------------------------------------------------

For modulus \(M\), base-6 row sector is defined by
\begin{equation}
    s_{6,M}(r)
    =
    \left\lfloor 6\cdot\frac{r\bmod M}{M}\right\rfloor
    \in\{0,1,2,3,4,5\},
    \label{eq:sector}
\end{equation}
where \(r\) is the row-window midpoint.

For each sector \(s\), let \(E_s\) be the number of expansion-tail events and \(C_s\) the number of compression-tail events. The sector imbalance vector is
\begin{equation}
    I_s=\log\left(\frac{E_s+1/2}{C_s+1/2}\right),
\end{equation}
and the sector-amplitude statistic is
\begin{equation}
    \Ssector(M)
    =
    \sqrt{
    \frac{1}{6}\sum_{s=0}^{5}
    \left(I_s-\overline{I}\right)^2
    }.
    \label{eq:ssector}
\end{equation}

The earlier preprint used local permutation guardrails after candidate moduli had been identified. Those guardrails showed that several selected moduli, especially \(1009\) and \(990\), appeared nontrivial under local tests. The missing step was a search-adjusted null over the modulus-selection process itself.

%--------------------------------------------------------------------
\section{Why the original positive claim is no longer maintained}
%--------------------------------------------------------------------

The original row-sector claim had three components:
\begin{enumerate}
    \item a sector-amplitude asymmetry exists in the terminal residual tail population;
    \item modulus \(1009\) is the most consistent tested sector-amplitude ruler;
    \item the row-sector structure is empirically orthogonal to the \(h_1\) packet corridor.
\end{enumerate}

The third component, orthogonality to \(h_1\), remains historically true in the narrow sense that the attempted \(h_1\)-coupling tests did not confirm \(h_1\) gating or local \(h_1\times\)sector modulation. However, the first two components no longer survive as positive claims. After the companion paper's \(h_1\) explanation was revised to altitude architecture, the row-sector signal required its own search-adjusted test. P86 supplies that test and rejects the positive modulus interpretation.

Thus the present paper is not a continuation of the old positive claim. It is its validation audit and replacement interpretation.

%--------------------------------------------------------------------
\section{Dense modulus search-adjusted validation}
%--------------------------------------------------------------------

\subsection{P86 audit design}

The P86 validation audit uses the following frozen design:
\begin{itemize}
    \item primary metric: \texttt{Gap\_to\_Fitted\_Ratio};
    \item windows: \(W\in\{900,1000,1100,1300\}\);
    \item tail threshold: \(|z|\ge 2.575\);
    \item modulus scan: \(M=800,\ldots,1200\);
    \item sectorization: base-6;
    \item permutations: \(500\) tail-label shuffles per configuration.
\end{itemize}

For each real configuration, \(S_{\mathrm{sector}}(M)\) is computed for all \(401\) moduli. For each permutation, tail signs are shuffled while row positions are held fixed, \(S_{\mathrm{sector}}(M)\) is recomputed for all \(401\) moduli, and the null retains
\[
\max_{800\le M\le 1200} S_{\mathrm{sector}}^{\mathrm{null}}(M).
\]
The search-adjusted \(p\)-value compares the real selected statistic to this null distribution of maximums.

\subsection{Primary dense-scan result}

Table~\ref{tab:p86-summary} gives the primary P86 summary. No tested configuration supports the earlier modulus-\(1009\) claim.

\begin{table}[H]
\centering
\small
\begin{adjustbox}{max width=\textwidth}
\begin{tabular}{rrrrrrrr}
\toprule
\textbf{\(W\)} & \textbf{Tail count} & \textbf{Best \(M\)} & \textbf{Best \(S\)} & \textbf{\(S_{1009}\)} & \textbf{Rank 1009} & \textbf{\(p_{\mathrm{best}}\)} & \textbf{\(p_{1009}\)} \\
\midrule
900  & 112 & 1133 & 1.250916 & 0.476623 & 48  & 0.079840 & 1.000000 \\
1000 & 104 & 1087 & 1.172635 & 0.385867 & 151 & 0.235529 & 1.000000 \\
1100 & 93  & 1133 & 0.764446 & 0.395511 & 143 & 1.000000 & 1.000000 \\
1300 & 71  & 1022 & 1.127564 & 0.591068 & 58  & 0.938124 & 1.000000 \\
\bottomrule
\end{tabular}
\end{adjustbox}
\caption{P86 dense modulus search-adjusted validation. The null scans \(M=800,\ldots,1200\) for every permutation and keeps its maximum \(S_{\mathrm{sector}}\). Neither \(M=1009\) nor the best real modulus survives in any tested configuration.}
\label{tab:p86-summary}
\end{table}

Across these configurations:
\[
\#\{W:\ p_{1009}\le 0.05\}=0,
\]
\[
\#\{W:\ p_{\mathrm{best}}\le 0.05\}=0,
\]
and
\[
\#\{W:\ \mathrm{rank}(1009)\le 10\}=0.
\]
The P86 status is therefore
\[
\texttt{P86\_NO\_MODULUS\_SURVIVES\_LOOK\_ELSEWHERE}.
\]

\optionalfigure{p86_dense_modulus_20260504_045931_primary_dense_scan.png}
{Primary dense modulus scan from P86. The real curve \(S_{\mathrm{sector}}(M)\) is plotted over \(M=800,\ldots,1200\), with markers for \(M=1009\), \(M=990\), and the best real modulus. The search-adjusted null ceiling is obtained from permutation replicates that scan the same modulus range.}
{fig:p86-dense-scan}

\optionalfigure{p86_dense_modulus_20260504_045931_primary_search_adjusted_null.png}
{P86 search-adjusted null. Each null replicate shuffles tail signs and records the maximum sector amplitude over \(M=800,\ldots,1200\). The real modulus-\(1009\) statistic lies well inside the null envelope.}
{fig:p86-null}

\subsection{Interpretation}

The dense null shows that the earlier sector result is compatible with modulus selection over sparse tail populations. When hundreds of moduli are searched, it is expected that some modulus will produce a large apparent \(S_{\mathrm{sector}}\) by chance. The correct null must search the same family, and under that null the signal disappears.

This is not a numerical bug or a failure of the sector statistic to compute what it was designed to compute. It is a post-selection failure: the local statistic was real conditional on a chosen modulus, but the chosen modulus was not protected against the look-elsewhere effect.

%--------------------------------------------------------------------
\section{Claim transition ledger}
%--------------------------------------------------------------------

Table~\ref{tab:claim-ledger} records the transition from the earlier row-sector preprint to the present replacement interpretation.

\begin{table}[H]
\centering
\small
\begin{adjustbox}{max width=\textwidth}
\begin{tabular}{p{0.30\textwidth}p{0.35\textwidth}p{0.28\textwidth}}
\toprule
\textbf{Original claim} & \textbf{Validation result} & \textbf{Replacement statement} \\
\midrule
A row-sector rebound structure exists as a positive discovery. &
The apparent sector amplitudes do not survive dense search-adjusted modulus controls. &
Rejected as a search-hardened positive claim. \\
\midrule
Modulus \(1009\) is the leading sector-amplitude ruler. &
In P86, \(M=1009\) survives in zero configurations and is never ranked in the top ten across the tested windows. &
Modulus \(1009\) is not supported as a physical or distinguished modulus. \\
\midrule
Modulus \(990\) is co-significant with \(1009\). &
Under dense look-elsewhere correction, the modulus family as a whole fails to produce a surviving positive result. &
The \(990/1009\) framing is retired. \\
\midrule
The sector structure is orthogonal to \(h_1\). &
The original \(h_1\) claim itself is replaced by altitude architecture. The old orthogonality question no longer carries a positive sector result. &
The episode becomes a methodological audit of projection artifacts. \\
\bottomrule
\end{tabular}
\end{adjustbox}
\caption{Transition from the earlier row-sector interpretation to the present search-adjusted interpretation.}
\label{tab:claim-ledger}
\end{table}

%--------------------------------------------------------------------
\section{Relation to terminal altitude architecture}
%--------------------------------------------------------------------

The revised companion manuscript identifies a terminal altitude architecture in highly compressed adjacent-spacing packets. That architecture explains the old \(h_1\) Arc80 corridor as an altitude projection. The present paper shows that the row-sector interpretation does not supply a second surviving mechanism.

The resulting unified picture is:
\begin{enumerate}
    \item the terminal compression packet field is real and altitude-structured;
    \item the continuous \(h_1\) projection reveals that altitude structure but is not its mechanism;
    \item the discrete row-sector modulus projection does not survive dense modulus look-elsewhere controls;
    \item the surviving positive object is altitude architecture, not phase locking or modulus sectorization.
\end{enumerate}

This reframing is important. It prevents the validation failure of \(h_1\) and \(M=1009\) from being misread as failure of the packet object itself. The packet object survives; the projected coordinate stories do not.

%--------------------------------------------------------------------
\section{Discussion}
%--------------------------------------------------------------------

The negative result in this paper is scientifically useful because it closes a specific loophole. The earlier row-sector analysis used multiple local nulls, bin-resolution checks, and \(h_1\)-interaction audits. Those tests were valuable, but they did not answer the decisive question: whether the selected modulus survives a null that searches the same modulus family.

P86 answers that question. It does not.

This result illustrates a general hazard in empirical number theory. Long arithmetic sequences are rich enough that many coordinate systems can be imposed on them. A modulus scan, a phase projection, or a geometric sectorization can appear structured if evaluated only after selection. The null must reproduce the selection process.

\subsection{What remains from the original continuation paper}

Several components remain useful:
\begin{itemize}
    \item the formal row-sector statistic \(S_{\mathrm{sector}}\);
    \item the interaction-residual framework for separating main effects from local \(h_1\)-sector interactions;
    \item the bin-resolution calibration logic that prevented overclaiming local \(h_1\times\)sector structure;
    \item the documented failure modes of a plausible exploratory coordinate system.
\end{itemize}
These are methodological contributions, not positive evidence for a sector engine.

\subsection{What does not remain}

The following claims should no longer be made:
\begin{itemize}
    \item modulus \(1009\) is a physical modulus;
    \item modulus \(1009\) is a search-hardened sector-amplitude leader;
    \item modulus \(990\) and \(1009\) define a validated scale-separated sector pair;
    \item the row-sector rebound is a confirmed positive structure independent of post-selection.
\end{itemize}

%--------------------------------------------------------------------
\section{Limitations}
%--------------------------------------------------------------------

The dense scan in P86 is intentionally focused on the primary fitted-ratio seam-scale family. It rejects the earlier positive claim under the relevant search-adjusted family \(M=800,\ldots,1200\), but it does not rule out every possible sectorization under every possible theoretical prior. A future theory that pre-specifies a modulus and a sector rule before seeing the data would require a different validation design.

The tail populations are sparse. This sparsity is part of why the look-elsewhere effect is severe: large apparent sector amplitudes can arise by chance when many moduli are scanned over a small number of tail events.

The result is finite-horizon. It concerns the loaded terminal LMFDB-derived range up to approximately \(30.610\mathrm{B}\). It does not claim that all future modulus-based analyses must fail.

%--------------------------------------------------------------------
\section{Conclusion}
%--------------------------------------------------------------------

The earlier row-sector preprint reported an apparent modulus-\(1009\) sector structure in terminal zeta-zero spacing residuals. A dense search-adjusted validation audit does not support that interpretation. When the null is allowed to scan the same \(M=800,\ldots,1200\) modulus family, neither modulus \(1009\) nor the best real modulus survives in any tested seam-scale configuration.

The row-sector result is therefore best understood as a look-elsewhere artifact produced by searching sparse tail-sign data over many moduli. In the unified replacement interpretation, the surviving empirical object is not a modulus sector engine but a terminal altitude architecture of highly compressed adjacent-spacing packets. The row-sector analysis remains valuable as a methodological audit: it shows how convincing a selected modulus can appear before search-adjusted validation is applied.

%--------------------------------------------------------------------
\section*{Data and Artifact Availability}
%--------------------------------------------------------------------

The original row-sector preprint was deposited under DOI
\href{https://doi.org/\paperdoi}{\paperdoi}. The principal companion preprint was deposited under DOI
\href{https://doi.org/\principaldoi}{\principaldoi}. The replacement validation chain generated outputs in
\[
\texttt{ForgeV16c\_P86\_DENSE\_MODULUS\_outputs}
\]
including the dense real modulus curve, permutation null distribution, configuration summary, and flat manifest. The broader replacement scaffold is recorded in
\[
\texttt{ForgeV16c\_P87\_PIVOT\_SCAFFOLD\_outputs}.
\]
The source zero ordinates are attributed to LMFDB and Odlyzko zero tables.

%--------------------------------------------------------------------
\section*{Use of AI Tools}
%--------------------------------------------------------------------

Generative AI tools were used for audit design review, code critique, adversarial claim evaluation, and draft assistance. All research design decisions, audit execution, interpretation, and final manuscript responsibility remain with Jeffery Huckstead.

%--------------------------------------------------------------------
\section*{Author Contributions}
%--------------------------------------------------------------------

Jeffery Huckstead conceived and executed the audit chain, designed the claim ledger and validation protocol, and wrote the manuscript.

%--------------------------------------------------------------------
\section*{Acknowledgments}
%--------------------------------------------------------------------

The author thanks the structured audit-review and adversarial critique process that shaped the ForgeV16c analytical chain.

%--------------------------------------------------------------------
\begin{thebibliography}{9}

\bibitem{huckstead_h1_principal}
J.~Huckstead,
\emph{Phase-Locked Extreme Compression in Adjacent Zeta-Zero Spacing:
  A Seam-Stitched Empirical Study Through the Full Available LMFDB
  Horizon},
Zenodo preprint, 2026.
DOI: \href{https://doi.org/\principaldoi}{\principaldoi}.

\bibitem{huckstead_altitude_replacement}
J.~Huckstead,
\emph{Terminal Altitude Architecture and Post-Selection Artifacts in Adjacent Zeta-Zero Spacing},
replacement manuscript in preparation, 2026.

\bibitem{riemann1859}
B.~Riemann,
\emph{Ueber die Anzahl der Primzahlen unter einer gegebenen Gr\"osse},
Monatsberichte der Berliner Akademie, 1859.

\bibitem{odlyzko_tables}
A.~M. Odlyzko,
\emph{Tables of zeros of the Riemann zeta function},
available at
\url{http://www.dtc.umn.edu/~odlyzko/zeta_tables/},
accessed 2026-05-03.

\bibitem{montgomery1973}
H.~L. Montgomery,
\emph{The pair correlation of zeros of the zeta function},
in \emph{Analytic Number Theory},
Proc.\ Sympos.\ Pure Math., vol.~24,
Amer.\ Math.\ Soc., Providence, RI, 1973, pp.~181--193.

\bibitem{mehta_random_matrices}
M.~L. Mehta,
\emph{Random Matrices},
3rd~ed., Elsevier, Amsterdam, 2004.

\end{thebibliography}

\end{document}